\properheading{Seventeenth Sunday in Ordinary Time}{Year A \textperiodcentered{} Lectionary no.~109 \textperiodcentered{} Ordinary Time Week XVII formulary}{\emph{Roman Missal}, Third Edition, and \emph{Lectionary for Mass}, Second Typical Edition, for the dioceses of the United States of America}

\begin{unitstable}
\unitphase{Introductory Rites}
\unitrow{Entrance Antiphon}{\emph{Cf.} Ps 67:6--7, 36 (Vulgate/Septuagint numbering; Ps 68 Hebrew). Latin incipit \incipit{Deus in loco sancto suo}. An approved entrance chant may be used instead.}
\unitrow{Penitential Act}{No proper text. On a Sunday the blessing and sprinkling of water may replace it; that rite has its own texts.}
\unitrow{Collect}{\incipit{Protéctor in te sperántium, Deus}. One text, shared by Years A, B and C.}
\unitphase{Liturgy of the Word}
\unitrow{First Reading}{1 Kings 3:5, 7--12. Verse 6 falls inside the span and is not proclaimed.}
\unitrow{Responsorial Psalm}{Ps 119:57, 72, 76--77, 127--128, 129--130; response \emph{Cf.} Ps 119:97a. Verses from six of the psalm's twenty-two alphabetic stanzas; the refrain from a seventh.}
\unitrow{Second Reading}{Romans 8:28--30.}
\unitrow{Gospel Acclamation}{\emph{Cf.} Matthew 11:25, set between the Alleluia.}
\unitrow{Gospel}{Matthew 13:44--52, \emph{or} the shorter form Matthew 13:44--46. Only one is proclaimed.}
\unitphase{Liturgy of the Eucharist}
\unitrow{Offertory chant}{No Missal antiphon. Any chant used is a local selection.}
\unitrow{Prayer over the Offerings}{\incipit{Súscipe, qu\'{\ae}sumus, Dómine, múnera}. Shared by Years A, B and C.}
\unitrow{Preface}{None proper. A Preface of Sundays in Ordinary Time is used unless the Eucharistic Prayer chosen carries its own.}
\unitrow{Eucharistic Prayer}{No proper insert. The choice is coupled to the Preface decision above.}
\unitphase{Communion Rites}
\unitrow{Communion Antiphon}{Either Ps 102:2 (Vulgate/Septuagint; Ps 103 Hebrew), \incipit{Bénedic, ánima mea, Dómino}; \emph{or} Matthew 5:7--8, \incipit{Beáti misericórdes}. The Missal prints them as a closed pair.}
\unitrow{Prayer after Communion}{\incipit{Súmpsimus, Dómine, divínum sacraméntum}. Shared by Years A, B and C.}
\end{unitstable}

\begin{foursenses}
Literal & A young king, newly succeeded and calling himself a child, is offered whatever he will ask and asks instead for a heart able to hear and to judge; the psalm's speaker prefers instruction to gold; Paul tells Roman believers that God works toward good for those who love him and has purposed their conformity to his Son; and Jesus, closing a discourse in parables, likens the kingdom to a buried treasure, one costly pearl, a net taking fish of every kind, and a trained scribe who produces new things and old.\\
Allegorical & The Latin tradition identified the hidden treasure variously as Christ concealed in the flesh (Hilary), as the sense hidden under the letter of Scripture (Origen, Jerome), as heavenly desire hidden in ascetical discipline (Gregory the Great), and as sacred doctrine within the Church (Aquinas, who reports the first three without adjudicating). These are competing figures, and the disagreement is reported rather than resolved.\\
Moral & The petitions turn value into conduct: the Collect asks that passing goods be so used now that what abides may already be held, and the Prayer over the Offerings returns what came from God's own bounty. Solomon's refusal of long life, riches and his enemies' lives is the reading's own contrast, and the psalm's ranking of the commandments above gold and topaz is the assembly's answer to it.\\
Anagogical & Only the longer Gospel form reaches the shore, the sorting and the furnace; the shorter ends in joy and purchase. Romans carries its chain of verbs through to \emph{glorified}, a past tense for what is not yet seen, and the Prayer after Communion names the sacrament a perpetual memorial of the Passion ordered to salvation. What is hoped for is adherence to what remains, not the securing of a possession.\\
\end{foursenses}
